% Created 2026-10-08 Thu 09:38
% Intended LaTeX compiler: pdflatex
\documentclass[11pt]{article}
\usepackage[utf8]{inputenc}
\usepackage[T1]{fontenc}
\usepackage{graphicx}
\usepackage{longtable}
\usepackage{wrapfig}
\usepackage{rotating}
\usepackage[normalem]{ulem}
\usepackage{amsmath}
\usepackage{amssymb}
\usepackage{capt-of}
\usepackage{hyperref}
\author{fcalado}
\date{\today}
\title{}
\hypersetup{
 pdfauthor={fcalado},
 pdftitle={},
 pdfkeywords={},
 pdfsubject={},
 pdfcreator={Emacs 29.3 (Org mode 9.6.15)}, 
 pdflang={English}}
\begin{document}

\tableofcontents

While this kind of critique looks for \textbf{\textbf{common ground}} between cis
and trans bodily experience, the representation of trans subjects is a
separate issue. As scholarship in media studies have argued, trans
studies subjects occupy distinct roles and thus require their own
models of analysis. Thomas Billard and Erique Zhang explain that
existing queer and feminist rubrics, which operate within (and
against) contexts of patriarchy and heteronormativity, do not apply in
the same way to the experience of trans subjects. They ask, for
example, "Where does the trans man fall in the hierarchy of
patriarchal domination? Does he now dominate women by virtue of his
transition?" (Billard and Zhang 195). The authors argue that trans
media subjects need their own frameworks for representation, one that
is based on reading how trans people represent themselves.

But, the major argument of this paper is that the experience of cis
subjects might be read under trans studies frameworks, albeit within
certain constraints. While cis and trans media representations vastly
contrast each other, with trans subjects occupying distinct roles and
plots from cis ones, they nonetheless may overlap temporarily. For
example, as Cáel M. Keegan points out, the line between cis and trans
dissolves where gender norms are experienced as enforced and
oppressive. He discusses the film \emph{The Assignment} (2016), where a
hitman named Frank is punished by being forceably transitioned into a
woman by a rogue surgeon. Keegan argues that Frank evokes a plight of
a trans man: "Frank is forced to live and behave as a woman because he
has been assigned a female body\ldots{} [the film] is arguably about the
experience of transgender men---a point entirely missed by the press
coverage" (Keegan 34). In Keegan's argument, such "bad objects" for
trans media representation, which represent trans lives through tropes
of suffering and stereotype, eventually lead to "a world in which the
distinction between cis and trans ceases to exist altogether, because
the systems enforcing binary sex and gender are dismantled" (36).
\end{document}
